\documentclass[11pt,letterpaper]{article}

% -- Fonts --
\usepackage{fontspec}
\setmainfont{TeX Gyre Termes}
\setsansfont{TeX Gyre Heros}
\setmonofont{TeX Gyre Cursor}

% -- Page geometry --
\usepackage[margin=1.15in]{geometry}

% -- Typographic refinement --
\usepackage{microtype}
\emergencystretch=2em

% -- Colors --
\usepackage{xcolor}
\definecolor{pscpurple}{HTML}{4A148C}
\definecolor{pscblue}{HTML}{1565C0}
\definecolor{pscgold}{HTML}{F57F17}
\definecolor{pscdark}{HTML}{263238}
\definecolor{psclight}{HTML}{F3E5F5}

% -- Tables --
\usepackage{booktabs}
\usepackage{tabularx}
\newcolumntype{Y}{>{\raggedright\arraybackslash}X}

% -- Headers/footers --
\usepackage{fancyhdr}
\pagestyle{fancy}
\fancyhf{}
\fancyhead[L]{\small\color{pscpurple}\textsc{The Infrastructure Dividend}}
\fancyhead[R]{\small\color{pscpurple}\textsc{PSC Research Series}}
\fancyfoot[C]{\small\thepage}
\renewcommand{\headrulewidth}{0.4pt}

% -- Section styling --
\usepackage{titlesec}
\titleformat{\section}{\Large\bfseries\color{pscpurple}}{\thesection}{1em}{}[\vspace{-0.5em}\rule{\textwidth}{0.4pt}]
\titleformat{\subsection}{\large\bfseries\color{pscblue}}{\thesubsection}{1em}{}

% -- Hyperlinks --
\usepackage{hyperref}
\hypersetup{colorlinks=true, linkcolor=pscblue, urlcolor=pscblue}

% -- Misc --
\usepackage{enumitem}
\usepackage{parskip}

\begin{document}

% ===========================================
% TITLE PAGE
% ===========================================
\begin{titlepage}
\centering
\vspace*{2cm}

{\small\textsc{\color{pscgold}PSC Research Series --- Political Science, History \& the Complex Plane}}

\vspace{0.5cm}
\rule{\textwidth}{1.5pt}
\vspace{0.8cm}

{\Huge\bfseries\color{pscpurple} The Infrastructure\\[0.3cm]
Dividend}

\vspace{0.6cm}
\rule{0.6\textwidth}{0.4pt}
\vspace{0.6cm}

{\Large\itshape UBI as Return on Collective Investment}

\vspace{1.5cm}

{\large\textsc{An Essay on Governance, Infrastructure, and the\\Accounting Correction That Reframes Everything}}

\vspace{0.5cm}

{\normalsize The infrastructure was built. The taxes were paid.\\
The returns are being captured privately.\\
The dividend corrects the ledger.}

\vfill

{\small
Date: April 7, 2026 \(\vert\) Classification: Research Essay\\
Branch: \texttt{artemis-ii} \(\vert\) Artemis II Mission --- Flight Day 7, Return Coast\\
Series: \texttt{tibsfox.com/Research/PSC/}
}

\vspace{0.5cm}
\rule{\textwidth}{1.5pt}
\vspace{0.3cm}
{\footnotesize Tibsfox Inc.\ \(\cdot\) Mukilteo, Washington}
\end{titlepage}

% ===========================================
\section{The Question Nobody Asks}
% ===========================================

Every corporation that moves a product relies on roads it did not build. Every bank that processes a transaction relies on a legal system it did not create. Every tech company that serves a customer relies on an electrical grid, a telecommunications network, a water system, a school system that educated its workforce, and a military that secures the territory in which it operates. These are not gifts. They are infrastructure---built over centuries, paid for by generations of taxpayers, maintained by public expenditure, and used for private profit.

The question is not whether a Universal Basic Income is affordable. The question is: why don't the people who paid for the infrastructure receive a dividend from the businesses that use it?

% ===========================================
\section{The Infrastructure Stack}
% ===========================================

Infrastructure is not just roads and bridges. It is a layered system, accumulated over centuries, each layer enabling the layers above it:

\subsection{Layer 1 --- Physical Infrastructure}
Roads, bridges, ports, airports, rail, waterways, dams, levees. The Interstate Highway System alone cost \$530 billion (2024 dollars) and took 35 years to build. Every logistics company, every commuter, every Amazon delivery uses it daily. The Army Corps of Engineers maintains 25,000 miles of navigable waterways. The FAA operates the airspace. None of this was built by the private sector.

\subsection{Layer 2 --- Utility Infrastructure}
The electrical grid (160,000+ miles of high-voltage transmission lines, largely built with public subsidies and eminent domain). Water systems. Sewage. Natural gas pipelines. The internet backbone---ARPANET was a Defense Department project; the fiber that carries global commerce was laid with public rights-of-way and taxpayer-funded research.

\subsection{Layer 3 --- Legal Infrastructure}
The court system. Contract enforcement. Property rights. Patent law. Corporate law itself---the limited liability corporation is a legal construct that socializes risk while privatizing profit. Without the legal system, no contract is enforceable, no property is secure, no corporation exists.

\subsection{Layer 4 --- Knowledge Infrastructure}
Public education. State universities. The research funding pipeline (NIH, NSF, DARPA, NASA). The pharmaceutical industry depends on NIH-funded basic research. The tech industry depends on computer science research funded by DARPA and NSF. GPS---built by the military, given to the world---underpins everything from agriculture to ride-sharing to financial trading.

\subsection{Layer 5 --- Social Infrastructure}
Healthcare systems. Social Security. Unemployment insurance. Food safety regulation. Environmental protection. The stability that allows commerce to function. Businesses do not thrive in failed states---they thrive in societies where workers are healthy, educated, fed, and not in armed conflict with each other.

\subsection{Layer 6 --- Security Infrastructure}
Military. Police. Fire. Emergency services. Diplomatic corps. Intelligence services. The entire apparatus that maintains territorial integrity and internal order. The cost is enormous---the United States spends over \$800 billion annually on defense alone---and the benefit accrues disproportionately to those with the most property and commerce to protect.

\medskip
\begin{center}
\colorbox{psclight}{\parbox{0.85\textwidth}{\centering\itshape Every business in America operates on top of all six layers simultaneously.\\None built any of them. All profit from them.}}
\end{center}

% ===========================================
\section{The Historical Precedent: Dividends on Shared Resources}
% ===========================================

This is not a radical idea. It is the oldest idea in governance.

\textbf{The Commons.} For most of human history, shared resources---fisheries, forests, grazing land, waterways---were managed collectively, and the yield was distributed. The enclosure movements in England (1604--1914) privatized the commons, converting shared wealth into private property. The dispossessed were told to find wages. The wages never caught up.

\textbf{Alaska Permanent Fund.} Since 1982, every resident of Alaska has received an annual dividend from the state's oil wealth---typically \$1,000--\$2,000 per year. The principle is simple: the oil belongs to the people of Alaska; companies pay royalties to extract it; a portion of those royalties is invested; the returns are distributed equally. This is an infrastructure dividend. It works. It has bipartisan support. It reduces poverty. And it is funded not by taxes but by resource extraction fees.

\textbf{Norway's Government Pension Fund.} Worth over \$1.7 trillion, built from North Sea oil revenues. Norway's approach: the resources belong to the nation; extraction profits are invested for the benefit of all citizens, present and future. The fund now owns approximately 1.5\% of all listed shares globally.

\textbf{Singapore's CPF and Public Housing.} Over 80\% of Singaporeans live in government-built housing. The Central Provident Fund is a mandatory savings and investment system. Singapore treated infrastructure investment as a national return-generating asset, not an expense.

The pattern is consistent: when nations treat shared resources as collective investments rather than political expenditures, the returns can be distributed. The infrastructure dividend is not charity. It is a return on investment.

% ===========================================
\section{The Math: GDP as Infrastructure Yield}
% ===========================================

United States GDP in 2025: approximately \$29 trillion.

Aschauer (1989) found that a 1\% increase in public capital stock produced a 0.39\% increase in private output. Fernald (1999) found that roads directly contributed to productivity growth. The Congressional Budget Office estimates the rate of return on public infrastructure investment at 15--30\%.

Conservative estimate: public infrastructure enables at least 15\% of GDP---approximately \$4.35 trillion annually.

\begin{center}
\begin{tabularx}{0.85\textwidth}{Y r r}
\toprule
\textbf{Dividend Rate} & \textbf{Total Pool} & \textbf{Per Person\slash Year} \\
\midrule
2\% of GDP & \$580 billion & \$2,230 \\
3\% of GDP & \$870 billion & \$3,346 \\
5\% of GDP & \$1.45 trillion & \$5,577 \\
\bottomrule
\end{tabularx}
\end{center}

\smallskip
Based on approximately 260 million American adults. For context: current U.S.\ federal revenue is approximately \$4.4 trillion. The infrastructure dividend at 2\% of GDP would represent 13\% of federal revenue.

Comparable existing expenditures:
\begin{itemize}[nosep,leftmargin=1.5em]
\item Defense: \$886 billion (2024)
\item Social Security: \$1.4 trillion
\item Medicare + Medicaid: \$1.6 trillion
\item Interest on national debt: \$659 billion
\end{itemize}

The question is not whether the money exists. It is whether the political will exists to redirect infrastructure returns to the people who funded the infrastructure.

% ===========================================
\section{The Global Frame}
% ===========================================

\textbf{European social democracies} already return a higher percentage of GDP to citizens through healthcare, education, transit, and social programs. The Scandinavian model---high taxes, high services, high quality of life---is effectively an infrastructure dividend delivered as services rather than cash.

\textbf{Developing nations} face a different version. When infrastructure is built with foreign loans, the returns flow to creditors rather than citizens. The debt trap---where nations borrow to build infrastructure whose profits are extracted by multinational corporations who then repatriate earnings to tax havens---is the infrastructure dividend flowing in reverse.

\textbf{Indigenous nations} understood this intuitively. The potlatch economies of the Pacific Northwest distributed surplus as a matter of governance. The commons were managed for collective benefit. Sovereignty meant that the people who lived on the land benefited from what the land produced.

This connects directly to the Ring of Fire sovereignty research: the cooperative infrastructure model (KNET, Wasaya Airways, Marten Falls equity co-investment) is the infrastructure dividend principle applied at the community level. When the Marten Falls First Nation takes an equity position in the Northern Road Link, they are claiming their infrastructure dividend---not as charity or consultation theater, but as co-owners of infrastructure that crosses their territory.

% ===========================================
\section{Why ``Dividend'' and Not ``Welfare''}
% ===========================================

Language matters.

A \textbf{welfare program} implies dependency. It implies that the recipient has failed and requires assistance. It carries stigma. It is politically vulnerable because it can be framed as ``giving money to people who don't work.''

A \textbf{dividend} implies ownership. It implies that the recipient has invested and is receiving a return. It carries respect. It is politically durable because it is framed as ``receiving your share of what you already paid for.''

The Alaska Permanent Fund is politically untouchable precisely because Alaskans understand it as \emph{their dividend}, not a handout. Every Alaskan, regardless of income, receives the same amount. The universality is the point.

The infrastructure dividend reframes UBI not as a progressive policy experiment but as a conservative accounting correction. The infrastructure was built. The taxes were paid. The returns are being captured privately. The dividend corrects the ledger.

% ===========================================
\section{The Corporate Counterargument and Why It Fails}
% ===========================================

The standard objection: corporations already pay taxes.

U.S.\ corporate tax revenue in 2023 was approximately \$420 billion---roughly 1.4\% of GDP. The effective corporate tax rate, after deductions, credits, and offshore structures, averages 12--15\% for large corporations, well below the statutory 21\%.

Meanwhile, the infrastructure those corporations use costs the public approximately \$1.2 trillion annually in direct infrastructure spending (federal + state + local).

\begin{center}
\colorbox{psclight}{\parbox{0.85\textwidth}{\centering Corporations are paying \$420 billion for access to \$1.2 trillion in infrastructure.\\That is a \textbf{65\% subsidy.}\\The infrastructure dividend reclaims a fraction of that gap.}}
\end{center}

% ===========================================
\section{The Complex Plane Reading}
% ===========================================

In the framework of this research series, the infrastructure dividend has both a real component and an imaginary component.

\textbf{\(R(t)\):} The institutional, measurable reality. GDP, tax revenue, infrastructure spending, poverty rates, Gini coefficients. These show that infrastructure enables wealth, wealth is concentrated, and the concentration is increasing.

\textbf{\(X(t)\):} The experiential, felt reality. The sense that the system is rigged. The anger when a billionaire pays a lower effective tax rate than a teacher. The erosion of trust in institutions. The ``deaths of despair'' documented by Case and Deaton. The populist backlash that drives democratic backsliding.

The infrastructure dividend addresses both axes simultaneously. On the real axis, it redistributes measurable wealth. On the imaginary axis, it restores the felt connection between contribution and return---the sense that the system recognizes your existence and your investment.

A system where \(R(t)\) is technically democratic but \(X(t)\) is experienced as extractive is a flawed democracy drifting toward the competitive authoritarianism attractor. The infrastructure dividend is a stabilizing force on both axes.

% ===========================================
\section{Implementation Principles}
% ===========================================

\begin{enumerate}[leftmargin=1.5em]
\item \textbf{Universal, not means-tested.} Like the Alaska Permanent Fund. Universality removes stigma, reduces bureaucracy, and builds political durability.

\item \textbf{Funded by infrastructure usage fees, not income taxes.} Carbon taxes, data extraction fees, spectrum licensing, financial transaction taxes, and resource extraction royalties are all infrastructure usage fees.

\item \textbf{Indexed to GDP, not legislated as a fixed amount.} When the economy grows, the dividend grows. When it contracts, it contracts. This aligns the dividend with the actual infrastructure yield.

\item \textbf{Layered by jurisdiction.} National infrastructure produces a national dividend. State infrastructure produces a state dividend. Tribal infrastructure produces a tribal dividend. Each layer of the stack has its own yield and its own constituency.

\item \textbf{Paired with infrastructure investment, not a substitute for it.} The dividend is a return on infrastructure. Cutting infrastructure to fund the dividend is eating the seed corn. The two must grow together.
\end{enumerate}

% ===========================================
\section*{Conclusion}
% ===========================================

The infrastructure dividend is not a new idea. It is the oldest idea in governance, expressed in modern economic terms. Every society that has ever managed shared resources has faced the question of how to distribute the returns. The ones that distributed equitably tended to be stable. The ones that concentrated returns tended to revolt.

The United States built the most extensive public infrastructure in human history---and then allowed the returns on that infrastructure to be captured almost entirely by private interests. The infrastructure dividend corrects this. Not as charity, not as welfare, not as a progressive policy experiment---but as a dividend. A return on an investment that every taxpayer, every generation, has already made.

\begin{quote}
\itshape The money is not missing. It is misallocated. The infrastructure exists. The returns are real. The only question is whether the people who built it will ever see their share.
\end{quote}

\vfill
\begin{center}
\rule{0.4\textwidth}{0.4pt}\\[0.5em]
{\small\color{pscpurple} Written during Artemis II Mission \(\cdot\) Flight Day 7, Return Coast \(\cdot\) April 2026\\
Itself a product of 60 years of publicly funded aerospace infrastructure---\\
from the Saturn V to the SLS, built by taxpayers,\\
used to push the boundaries of what humanity can do.\\[0.3em]
PSC Research Series \(\cdot\) \texttt{tibsfox.com/Research/PSC/}}
\end{center}

\end{document}
